% Generated from pinned audit evidence.
\newcommand{\TrainSeconds}{361.6}
\newcommand{\TrainMemory}{0.543}
\newcommand{\FirstIdentity}{1.17}
\newcommand{\FirstShuffled}{1.17}
\newcommand{\FirstZero}{0.00}
\newcommand{\BlockIdentity}{0.586}
\newcommand{\BlockGain}{0.098}
\newcommand{\BlockLow}{0.000}
\newcommand{\BlockHigh}{0.307}
\newcommand{\BlockFone}{0.320}
\newcommand{\RunDecision}{reject}

% Generated from the pinned V54 audit; no fitted presentation values.
\newcommand{\DensitySeconds}{207.8}
\newcommand{\DensityMemory}{0.400}
\newcommand{\DensityNll}{15.219}
\newcommand{\DensityShuffledNll}{15.254}
\newcommand{\DensityAdvantage}{0.0355}
\newcommand{\DensityLow}{-0.1014}
\newcommand{\DensityHigh}{0.1693}
\newcommand{\DensityIdentity}{1.562}
\newcommand{\DensityFone}{0.350}
\newcommand{\DensityCopyFone}{0.99949}
\newcommand{\DensityDecision}{reject}

% Generated from the locked V55 matched audit.
\newcommand{\JointNll}{7.026}
\newcommand{\ControlNll}{7.579}
\newcommand{\JointNllGain}{0.553}
\newcommand{\JointNllLow}{0.450}
\newcommand{\JointNllHigh}{0.665}
\newcommand{\JointTime}{342.4}
\newcommand{\ControlTime}{218.0}
\newcommand{\JointMemory}{0.648}
\newcommand{\ControlMemory}{0.400}
\newcommand{\JointSampleFone}{0.362}
\newcommand{\JointPrefixGain}{0.185}
\newcommand{\JointDecision}{reject}

\newcommand{\PaperTitle}{Learning conditional language\\in continuous image space}
\newcommand{\PaperSummary}{Written language includes forms that cannot be faithfully represented by a fixed character inventory. We study language modelling with writing images as both input and output, continuous internal states and one consumer GPU. A causal Chinese raster reader ranks the next glyph at 19.97\% top-1 on 2,048 development windows, above unigram, bigram and shuffled-history controls. A normalized autoregressive raster density preserves all 256 supplied development images under greedy copying. Joint reader--density training then improves natural-target likelihood over a matched frozen-reader continuation: \JointNll\ versus \ControlNll\ nats per image, a paired gain of \JointNllGain\ with 95\% source-record interval [\JointNllLow, \JointNllHigh]. Ordered earlier history supplies a further \JointPrefixGain{}-nat advantage over shuffling. The 18.71M-parameter joint runtime completes 8,000 additional updates in \JointTime\ seconds, peaking at \JointMemory\ GiB allocated CUDA memory, excluding foundation training. Autonomous target identity remains only 9/256, versus 8/256 with earlier history shuffled, and the registered generation criteria fail. These results demonstrate reliable visual copying and learned conditional image likelihood on one GPU, while leaving coherent language generation and broad instruction following unresolved.}
